\documentclass[10pt,letterpaper]{article}
\usepackage[margin=0.65in]{geometry}
\usepackage{amsmath}
\usepackage{url}
\usepackage{hyperref}
\usepackage{enumitem}
\usepackage{titlesec}
\usepackage{xcolor}

\definecolor{primary}{RGB}{15, 23, 42}
\definecolor{accent}{RGB}{2, 132, 199}

\hypersetup{
    colorlinks=true,
    linkcolor=accent,
    urlcolor=accent
}

\titleformat{\section}{\large\bfseries\color{primary}}{}{0em}{}[\titlerule]
\titlespacing*{\section}{0pt}{10pt}{4pt}

\setlist[itemize]{leftmargin=*,noitemsep,topsep=1pt,parsep=0pt}
\setlist[enumerate]{leftmargin=*,noitemsep,topsep=1pt,parsep=0pt}

\begin{document}
\pagestyle{empty}

\begin{center}
    {\huge \bfseries Rachit Bansal} \\[3pt]
    {\large \itshape Doctorate Student, Harvard University} \\[4pt]
    \href{https://rachitbansal.github.io}{rachitbansal.github.io} \ \textperiodcentered \ \href{mailto:rachitbansal@g.harvard.edu}{rachitbansal@g.harvard.edu} \ \textperiodcentered \ \href{https://scholar.google.com}{Google Scholar}
\end{center}

\vspace{-4pt}

\section*{Education}
\textbf{Harvard University}, \textit{Cambridge} \hfill \textbf{08/2024 -- Present} \\
\textit{Ph.D. in Computer Science (ongoing)}
\begin{itemize}
    \item Advised by Prof. Sham Kakade and Prof. David Alvarez-Melis.
    \item Part of the Harvard ML Foundations research group.
    \item Kempner Institute Graduate Fellow.
\end{itemize}

\vspace{3pt}

\textbf{Delhi Technological University}, \textit{India} \hfill \textbf{08/2018 -- 07/2022} \\
\textit{B.Tech. in Electrical Engineering}
\begin{itemize}
    \item Bachelor's thesis at the Technion, Israel (02/2022 -- 07/2022).
\end{itemize}

\section*{Experience}
\textbf{Meta}, \textit{California} \hfill \textbf{05/2025 -- 10/2025} \\
\textit{Research Scientist Intern at Meta Super-intelligence Labs (MSL)} with Aston Zhang

\vspace{3pt}

\textbf{Google DeepMind}, \textit{India} \hfill \textbf{07/2022 -- 07/2024} \\
\textit{Pre-doctoral Researcher} with Partha Talukdar and Prateek Jain

\vspace{3pt}

\textbf{Technion}, \textit{Israel} \hfill \textbf{09/2021 -- 07/2022} \\
\textit{Research Intern (Bachelor's Thesis)} with Yonatan Belinkov

\vspace{3pt}

\textbf{Adobe Research}, \textit{India} \hfill \textbf{01/2021 -- 09/2021} \\
\textit{Research Intern} with Balaji Krishnamurthy

\vspace{3pt}

\textbf{Google Summer of Code}, \textit{Remote | University of Oxford} \hfill \textbf{05/2020 -- 01/2021} \\
\textit{Contributor at CDLI} with Jacob Dahl

\section*{Publications}
\begin{enumerate}[label={[\arabic*]}]
    \item \textbf{Let's (not) just put things in Context: Test-Time Training for Long-Context LLMs} \\
    \underline{Rachit Bansal}, Aston Zhang, Rishabh Tiwari, Lovish Madaan, Sai Surya Duvvuri, Devvrit Khatri, David Brandfonbrener, David Alvarez-Melis, Prajjwal Bhargava, Mihir Sanjay Kale, Samy Jelassi \\
    \textit{International Conference on Learning Representations} \hfill [ \textbf{ICLR 2026} ]

    \item \textbf{The Art of Scaling Reinforcement Learning Compute for LLMs} \\
    Devvrit Khatri, Lovish Madaan, Rishabh Tiwari, \underline{Rachit Bansal}, Sai Surya Duvvuri, Manzil Zaheer, Inderjit S. Dhillon, David Brandfonbrener, Rishabh Agarwal \\
    \textit{International Conference on Learning Representations (Oral Presentation)} \hfill [ \textbf{ICLR 2026} ]

    \item \textbf{Adam or Gauss-Newton? A Comparative Study In Terms of Basis Alignment and SGD Noise} \\
    Bingbin Liu, \underline{Rachit Bansal}, Depen Morwani, Nikhil Vyas, David Alvarez-Melis, Sham M. Kakade \\
    \textit{International Conference for Machine Learning (under submission)} \hfill [ \textbf{ICML 2026} ]

    \item \textbf{LLM Augmented LLMs: Expanding Capabilities through Composition} \\
    \underline{Rachit Bansal}, Bidisha Samanta, Siddharth Dalmia, Nitish Gupta, Shikhar Vashishth, Sriram Ganapathy, Abhishek Bapna, Prateek Jain, Partha Talukdar \\
    \textit{International Conference on Learning Representations} \hfill [ \textbf{ICLR 2024} ]

    \item \textbf{Linear Connectivity Reveals Generalization Strategies} \\
    Jeevesh Juneja, \underline{Rachit Bansal}, Kyunghyun Cho, Jo\~ao Sedoc, Naomi Saphra \\
    \textit{International Conference on Learning Representations} \hfill [ \textbf{ICLR 2023} ]

    \item \textbf{Measures of Information Reflect Memorization Patterns} \\
    \underline{Rachit Bansal}, Danish Pruthi, Yonatan Belinkov \\
    \textit{Conference on Neural Information Processing Systems} \hfill [ \textbf{NeurIPS 2022} ]

    \item \textbf{Evaluating Explanations: How much do explanations from the teacher aid students?} \\
    Danish Pruthi, \underline{Rachit Bansal}, Bhuvan Dhingra, Livio Baldini Soares, Michael Collins, Zachary C. Lipton, Graham Neubig, William W. Cohen \\
    \textit{Transactions of the Association for Computational Linguistics} \\
    \textit{Presented at the Annual Conference for the Association of Computation Linguistics} \hfill [ \textbf{TACL 2022} ]

    \item \textbf{CoSe-Co: Text Conditioned Generative CommonSense Contextualizer} \\
    \underline{Rachit Bansal}, Milan Aggarwal, Sumit Bhatia, Jivat Kaur, Balaji Krishnamurthy \\
    \textit{North American Chapter of the Association for Computational Linguistics} \hfill [ \textbf{NAACL 2022} ]

    \item \textbf{LM-CORE: Language Models with Contextually Relevant External Knowledge} \\
    Jivat Kaur, Sumit Bhatia, Milan Aggarwal, \underline{Rachit Bansal}, Balaji Krishnamurthy \\
    \textit{North American Chapter of the Association for Computational Linguistics (Findings)} \hfill [ \textbf{NAACL 2022} ]

    \item \textbf{How Low is Too Low? A Computational Perspective on Extremely Low-Resource Languages} \\
    \underline{Rachit Bansal}, Himanshu Choudhary, Ravneet Punia, Niko Schenk, Jacob L Dahl, \'Emilie Pag\'e-Perron \\
    \textit{Student Research Workshop (SRW) at ACL} \hfill [ \textbf{ACL SRW 2021} ]
\end{enumerate}

\section*{Featured Academic Projects and Collaborations}
\textbf{Augmenting New Knowledge in Language Models through Composition} \hfill \textbf{07/2022 -- Present} \\
\textit{w/ Partha Talukdar, Prateek Jain, Nitish Gupta, Sid Dalmia} \hfill \textcolor{accent}{\textbf{Google Research}}
\begin{itemize}
    \item Worked as a part of a massive moonshot effort to create inclusive and equitable language representations.
    \item Led a large collaboration to introduce composition of language models as a paradigm to augment new knowledge.
    \item Proposed CALM: Using knowledge-specific models to augment new capabilities in a frozen language model. [ICLR'24]
    \item Working with Google DeepMind and the Bard team to test CALM for serving custom models to users.
\end{itemize}

\vspace{3pt}

\textbf{Relationship between Information Distribution and Model Behavior} \hfill \textbf{01/2022 -- 07/2022} \\
\textit{w/ Yonatan Belinkov, Danish Pruthi} \hfill \textcolor{accent}{\textbf{Technion}}
\begin{itemize}
    \item Evaluating generalization of neural models is difficult: Requires creation of labeled out-of-distribution sets.
    \item Employed information-theoretic metrics to study the information distribution across neurons as an intrinsic metric.
    \item For the first time, showed that such intrinsic metrics strongly correlate with generalization behaviors of a model.
    \item Demonstrated the usefulness of the study for model selection. [NeurIPS'22]
\end{itemize}

\vspace{3pt}

\textbf{Mode Connectivity in Loss Surfaces for Text Models} \hfill \textbf{10/2021 -- 10/2022} \\
\textit{w/ Naomi Saphra, Jo\~ao Sedoc, Kyunghyun Cho} \hfill \textcolor{accent}{\textbf{New York University}}
\begin{itemize}
    \item Analyzed linear model connectivity for multiple fine-tuned models from the same pre-trained language model.
    \item For the first time, observed clusters of models that lie in separate basins within the loss surface.
    \item Further observed that models belonging in the same cluster show identical generalization behaviors. [ICLR'23]
    \item Future work has utilized insights from our work for weight averaging and mechanistic interpretability.
\end{itemize}

\vspace{3pt}

\textbf{Teacher-Student Paradigm to Evaluate Model Explanations} \hfill \textbf{09/2020 -- 12/2021} \\
\textit{w/ Danish Pruthi, Bhuvan Dhingra, Zachary Lipton, Graham Neubig} \hfill \textcolor{accent}{\textbf{Carnegie Mellon University}}
\begin{itemize}
    \item A number of model explainability approaches exist but no means to quantitatively evaluate and measure progress.
    \item Established a student-teacher communication paradigm for automatic evaluation of explanations. [TACL'22]
\end{itemize}

\vspace{3pt}

\textbf{Grounding Language Models in Factual and Commonsense Knowledge} \hfill \textbf{01/2021 -- 09/2021} \\
\textit{w/ Milan Aggarwal, Sumit Bhatia, Balaji Krishnamurthy} \hfill \textcolor{accent}{\textbf{Adobe Research}}
\begin{itemize}
    \item Developed a framework to augment language model inputs with factual and commonsense knowledge on the fly.
    \item Demonstrated that our generic and efficient framework outperform large task-tuned models. [NAACL'22]
\end{itemize}

\vspace{3pt}

\textbf{Neural Machine Translation for Sumerian} \hfill \textbf{05/2020 -- 01/2021} \\
\textit{w/ Jacob Dahl, \'Emilie Pag\'e-Perron, Niko Schenk} \hfill \textcolor{accent}{\textbf{University of Oxford}}
\begin{itemize}
    \item Sumerian is the earliest written language in Mesopotamia and perhaps the world---dating back to 4th millennium BC.
    \item Led this open-source initiative with CDLI to adapt modern NMT for extremely low-resource languages [SRW, ACL'21].
    \item Built an end-to-end information extraction pipeline for Sumerian widely used by Sumerian assyriologists today.
\end{itemize}

\section*{Teaching and Featured Positions}
\textbf{Google Summer of Code}, Cuneiform Digital Library Initiative (CDLI). \textit{Mentor} \hfill \textbf{Summer 2022} \\
\textbf{Reinforcement Learning}, Coding Blocks. \textit{Student Instructor w/ Prateek Narang} \hfill \textbf{2020}
\begin{itemize}
    \item Recorded 10-hours worth of lectures and held a number of live webinars. Collaborated with course mentors to build project ideas, assignments, and quizzes.
\end{itemize}
\textbf{Foundations of Machine Learning \& Deep Learning}, Coding Blocks. \textit{Teaching Assistant w/ Prateek Narang} \hfill \textbf{2019}
\begin{itemize}
    \item Conducted classes and doubt sessions for a batch of 60 senior undergraduate students from all across the country. Built course quizzes and programming assignments in collaboration with other TAs.
\end{itemize}
\textbf{Reviewer}: ICLR'24, NeurIPS'23, EMNLP'23, ACL'23, NeurIPS'22

\section*{Featured Coursework}
\begin{itemize}
    \item \textbf{Mathematics}: Advanced Linear Algebra (2\textsuperscript{nd} Sem., DTU; \textit{University Rank-1}); MIT RES-6-012: Introduction to Probability, MIT OCW; Abstract Algebra, Group Theory, and Linear Algebra, IIT-KGP (NPTEL); Numerical and Engineering Optimization Methods (3\textsuperscript{rd} Sem., DTU); Swarm and Evolutionary Optimization (7\textsuperscript{th} Sem., DTU)
    \item \textbf{Machine Learning}: IFT 6760A: Matrix and tensor factorization techniques for machine learning, University of Montreal; MIT 18-065: Matrix Methods in Signal Processing, and Machine Learning, MIT OCW; Probabilistic Graphical Models Specialization, Stanford University; Bayesian Methods for Machine Learning, National Research University of Russia
    \item \textbf{Natural Language Processing}: CS11-737: Multilingual NLP, CMU; CS11-747: Neural Networks for NLP, CMU; Natural Language Processing (6\textsuperscript{th} Sem., DTU)
\end{itemize}

\end{document}
